% Options for packages loaded elsewhere
\PassOptionsToPackage{unicode}{hyperref}
\PassOptionsToPackage{hyphens}{url}
%
\documentclass[
  11pt,
]{article}
\usepackage{amsmath,amssymb}
\usepackage{iftex}
\ifPDFTeX
  \usepackage[T1]{fontenc}
  \usepackage[utf8]{inputenc}
  \usepackage{textcomp} % provide euro and other symbols
\else % if luatex or xetex
  \usepackage{unicode-math} % this also loads fontspec
  \defaultfontfeatures{Scale=MatchLowercase}
  \defaultfontfeatures[\rmfamily]{Ligatures=TeX,Scale=1}
\fi
\usepackage{lmodern}
\ifPDFTeX\else
  % xetex/luatex font selection
\fi
% Use upquote if available, for straight quotes in verbatim environments
\IfFileExists{upquote.sty}{\usepackage{upquote}}{}
\IfFileExists{microtype.sty}{% use microtype if available
  \usepackage[]{microtype}
  \UseMicrotypeSet[protrusion]{basicmath} % disable protrusion for tt fonts
}{}
\makeatletter
\@ifundefined{KOMAClassName}{% if non-KOMA class
  \IfFileExists{parskip.sty}{%
    \usepackage{parskip}
  }{% else
    \setlength{\parindent}{0pt}
    \setlength{\parskip}{6pt plus 2pt minus 1pt}}
}{% if KOMA class
  \KOMAoptions{parskip=half}}
\makeatother
\usepackage{xcolor}
\usepackage[margin=2.2cm]{geometry}
\usepackage{longtable,booktabs,array}
\usepackage{calc} % for calculating minipage widths
% Correct order of tables after \paragraph or \subparagraph
\usepackage{etoolbox}
\makeatletter
\patchcmd\longtable{\par}{\if@noskipsec\mbox{}\fi\par}{}{}
\makeatother
% Allow footnotes in longtable head/foot
\IfFileExists{footnotehyper.sty}{\usepackage{footnotehyper}}{\usepackage{footnote}}
\makesavenoteenv{longtable}
\setlength{\emergencystretch}{3em} % prevent overfull lines
\providecommand{\tightlist}{%
  \setlength{\itemsep}{0pt}\setlength{\parskip}{0pt}}
\setcounter{secnumdepth}{-\maxdimen} % remove section numbering
\ifLuaTeX
  \usepackage{selnolig}  % disable illegal ligatures
\fi
\IfFileExists{bookmark.sty}{\usepackage{bookmark}}{\usepackage{hyperref}}
\IfFileExists{xurl.sty}{\usepackage{xurl}}{} % add URL line breaks if available
\urlstyle{same}
\hypersetup{
  pdftitle={Worksheet 5C: Graphs of Sine and Cosine},
  hidelinks,
  pdfcreator={LaTeX via pandoc}}

\title{Worksheet 5C: Graphs of Sine and Cosine}
\author{}
\date{}

\begin{document}
\maketitle

\begin{quote}
\textbf{What This Worksheet Is About}

Sine and cosine don't just solve triangles --- they describe anything
that oscillates, cycles, or waves. Sound, light, electricity, ocean
tides, seasonal temperatures, and even the vibration of atoms all follow
sinusoidal patterns.

Understanding how to read and manipulate these graphs is essential for
signal processing, communications, and the Fourier analysis that
underlies much of modern technology.
\end{quote}

\begin{center}\rule{0.5\linewidth}{0.5pt}\end{center}

\hypertarget{part-a-the-basic-sine-and-cosine-curves-16-problems}{%
\subsection{Part A: The Basic Sine and Cosine Curves (16
problems)}\label{part-a-the-basic-sine-and-cosine-curves-16-problems}}

\hypertarget{the-parent-functions}{%
\subsubsection{The Parent Functions}\label{the-parent-functions}}

\[y = \sin(x) \quad \text{and} \quad y = \cos(x)\]

\textbf{Key Properties:} - \textbf{Period:} \(2\pi\) (one complete
cycle) - \textbf{Amplitude:} 1 (distance from center to peak) -
\textbf{Range:} \([-1, 1]\) - \textbf{Domain:} All real numbers

\textbf{The only difference:} Cosine is sine shifted left by
\(\frac{\pi}{2}\): \[\cos(x) = \sin\left(x + \frac{\pi}{2}\right)\]

\begin{center}\rule{0.5\linewidth}{0.5pt}\end{center}

\textbf{1.} Complete the table for \(y = \sin(x)\):

\begin{longtable}[]{@{}
  >{\raggedright\arraybackslash}p{(\columnwidth - 16\tabcolsep) * \real{0.0450}}
  >{\raggedright\arraybackslash}p{(\columnwidth - 16\tabcolsep) * \real{0.0450}}
  >{\raggedright\arraybackslash}p{(\columnwidth - 16\tabcolsep) * \real{0.1532}}
  >{\raggedright\arraybackslash}p{(\columnwidth - 16\tabcolsep) * \real{0.1532}}
  >{\raggedright\arraybackslash}p{(\columnwidth - 16\tabcolsep) * \real{0.1532}}
  >{\raggedright\arraybackslash}p{(\columnwidth - 16\tabcolsep) * \real{0.1532}}
  >{\raggedright\arraybackslash}p{(\columnwidth - 16\tabcolsep) * \real{0.0631}}
  >{\raggedright\arraybackslash}p{(\columnwidth - 16\tabcolsep) * \real{0.1622}}
  >{\raggedright\arraybackslash}p{(\columnwidth - 16\tabcolsep) * \real{0.0721}}@{}}
\toprule\noalign{}
\begin{minipage}[b]{\linewidth}\raggedright
\(x\)
\end{minipage} & \begin{minipage}[b]{\linewidth}\raggedright
\(0\)
\end{minipage} & \begin{minipage}[b]{\linewidth}\raggedright
\(\frac{\pi}{6}\)
\end{minipage} & \begin{minipage}[b]{\linewidth}\raggedright
\(\frac{\pi}{4}\)
\end{minipage} & \begin{minipage}[b]{\linewidth}\raggedright
\(\frac{\pi}{3}\)
\end{minipage} & \begin{minipage}[b]{\linewidth}\raggedright
\(\frac{\pi}{2}\)
\end{minipage} & \begin{minipage}[b]{\linewidth}\raggedright
\(\pi\)
\end{minipage} & \begin{minipage}[b]{\linewidth}\raggedright
\(\frac{3\pi}{2}\)
\end{minipage} & \begin{minipage}[b]{\linewidth}\raggedright
\(2\pi\)
\end{minipage} \\
\midrule\noalign{}
\endhead
\bottomrule\noalign{}
\endlastfoot
\(\sin(x)\) & & & & & & & & \\
\end{longtable}

\textbf{2.} Complete the table for \(y = \cos(x)\):

\begin{longtable}[]{@{}
  >{\raggedright\arraybackslash}p{(\columnwidth - 16\tabcolsep) * \real{0.0450}}
  >{\raggedright\arraybackslash}p{(\columnwidth - 16\tabcolsep) * \real{0.0450}}
  >{\raggedright\arraybackslash}p{(\columnwidth - 16\tabcolsep) * \real{0.1532}}
  >{\raggedright\arraybackslash}p{(\columnwidth - 16\tabcolsep) * \real{0.1532}}
  >{\raggedright\arraybackslash}p{(\columnwidth - 16\tabcolsep) * \real{0.1532}}
  >{\raggedright\arraybackslash}p{(\columnwidth - 16\tabcolsep) * \real{0.1532}}
  >{\raggedright\arraybackslash}p{(\columnwidth - 16\tabcolsep) * \real{0.0631}}
  >{\raggedright\arraybackslash}p{(\columnwidth - 16\tabcolsep) * \real{0.1622}}
  >{\raggedright\arraybackslash}p{(\columnwidth - 16\tabcolsep) * \real{0.0721}}@{}}
\toprule\noalign{}
\begin{minipage}[b]{\linewidth}\raggedright
\(x\)
\end{minipage} & \begin{minipage}[b]{\linewidth}\raggedright
\(0\)
\end{minipage} & \begin{minipage}[b]{\linewidth}\raggedright
\(\frac{\pi}{6}\)
\end{minipage} & \begin{minipage}[b]{\linewidth}\raggedright
\(\frac{\pi}{4}\)
\end{minipage} & \begin{minipage}[b]{\linewidth}\raggedright
\(\frac{\pi}{3}\)
\end{minipage} & \begin{minipage}[b]{\linewidth}\raggedright
\(\frac{\pi}{2}\)
\end{minipage} & \begin{minipage}[b]{\linewidth}\raggedright
\(\pi\)
\end{minipage} & \begin{minipage}[b]{\linewidth}\raggedright
\(\frac{3\pi}{2}\)
\end{minipage} & \begin{minipage}[b]{\linewidth}\raggedright
\(2\pi\)
\end{minipage} \\
\midrule\noalign{}
\endhead
\bottomrule\noalign{}
\endlastfoot
\(\cos(x)\) & & & & & & & & \\
\end{longtable}

\textbf{3.} Sketch both functions on the same axes from \(x = 0\) to
\(x = 2\pi\).

\textbf{4.} At what x-values does \(\sin(x) = \cos(x)\) in the interval
\([0, 2\pi]\)?

\textbf{5.} At what x-values does \(\sin(x) = 0\) in the interval
\([0, 2\pi]\)?

\textbf{6.} At what x-values does \(\cos(x) = 0\) in the interval
\([0, 2\pi]\)?

\textbf{7.} At what x-value does \(\sin(x)\) reach its maximum in
\([0, 2\pi]\)? Its minimum?

\textbf{8.} At what x-value does \(\cos(x)\) reach its maximum in
\([0, 2\pi]\)? Its minimum?

\begin{center}\rule{0.5\linewidth}{0.5pt}\end{center}

\hypertarget{symmetry}{%
\subsubsection{Symmetry}\label{symmetry}}

\textbf{9.} \(\sin(-x) ={}\) \_\_\_\_\_\_\_\_\_\_ (Sine is an
\textbf{odd} function)

\textbf{10.} \(\cos(-x) ={}\) \_\_\_\_\_\_\_\_\_\_ (Cosine is an
\textbf{even} function)

\textbf{11.} Verify: \(\sin(-\frac{\pi}{4}) = -\sin(\frac{\pi}{4})\)

\textbf{12.} Verify: \(\cos(-\frac{\pi}{3}) = \cos(\frac{\pi}{3})\)

\begin{center}\rule{0.5\linewidth}{0.5pt}\end{center}

** Science Connection: Sound Waves**

\textbf{13.} A pure musical tone is a sine wave. The note ``A'' above
middle C has a frequency of 440 Hz (cycles per second). - (a) What is
the period of this sound wave in seconds? - (b) How many complete cycles
occur in 1 second?

\textbf{14.} The equation for this sound wave (simplified) is:
\[y(t) = A\sin(2\pi \cdot 440 \cdot t)\] where \(t\) is time in seconds
and \(A\) is amplitude.

At \(t = 0\), what is \(y(0)\)?

\textbf{15.} At \(t = \frac{1}{880}\) seconds, what is \(y(t)\)?

\textbf{16.} At what time does the first peak occur?

\begin{center}\rule{0.5\linewidth}{0.5pt}\end{center}

\hypertarget{part-b-amplitude-vertical-stretch-14-problems}{%
\subsection{Part B: Amplitude --- Vertical Stretch (14
problems)}\label{part-b-amplitude-vertical-stretch-14-problems}}

\hypertarget{the-general-form}{%
\subsubsection{The General Form}\label{the-general-form}}

\[y = A\sin(x) \quad \text{or} \quad y = A\cos(x)\]

\textbf{Amplitude} \(= |A|\) (the distance from the midline to the peak)

\begin{itemize}
\tightlist
\item
  If \(A > 1\): vertical stretch (taller waves)
\item
  If \(0 < A < 1\): vertical compression (shorter waves)
\item
  If \(A < 0\): reflection across the x-axis
\end{itemize}

\begin{center}\rule{0.5\linewidth}{0.5pt}\end{center}

\textbf{Identify the amplitude and sketch one period:}

\textbf{17.} \(y = 3\sin(x)\)

\textbf{18.} \(y = 0.5\cos(x)\)

\textbf{19.} \(y = -2\sin(x)\)

\textbf{20.} \(y = -\cos(x)\)

\begin{center}\rule{0.5\linewidth}{0.5pt}\end{center}

\textbf{Write the equation of a sine function with:}

\textbf{21.} Amplitude 4, no other transformations.

\textbf{22.} Amplitude \(\frac{1}{3}\), reflected across the x-axis.

\begin{center}\rule{0.5\linewidth}{0.5pt}\end{center}

** Science Connection: Sound Intensity**

\textbf{23.} A quiet whisper might have amplitude 0.01 (in arbitrary
units), while a loud shout might have amplitude 1.0.

If quiet sound is \(y = 0.01\sin(x)\) and loud sound is
\(y = 1.0\sin(x)\): - (a) How many times larger is the amplitude of the
shout? - (b) Sound intensity is proportional to amplitude squared. How
many times more intense is the shout?

\textbf{24.} Noise-canceling headphones work by generating a sound wave
that is the \textbf{negative} of the incoming noise.

If noise is \(y = A\sin(x)\), what wave cancels it?

\begin{center}\rule{0.5\linewidth}{0.5pt}\end{center}

** Science Connection: Light Waves**

\textbf{25.} Light is an electromagnetic wave. The electric field
oscillates as a sine wave.

If blue light has a smaller amplitude than red light (same intensity),
which has: - (a) More energy per photon? (Hint: This depends on
frequency, not amplitude) - (b) More photons for the same total energy?

\textbf{26.} In interference experiments, two waves can: - Add
constructively (peaks align): \(\sin(x) + \sin(x) = 2\sin(x)\) - Add
destructively (peaks align with troughs):
\(\sin(x) + \sin(x + \pi) = 0\)

Verify the destructive case:
\(\sin(x) + \sin(x + \pi) = \sin(x) + (-\sin(x)) = 0\)

\begin{center}\rule{0.5\linewidth}{0.5pt}\end{center}

** Science Connection: Radio Waves**

\textbf{27.} AM radio (Amplitude Modulation) encodes information by
varying the amplitude of a carrier wave.

The carrier might be \(y = \sin(1000000 \cdot 2\pi t)\) (1 MHz).

When you speak, the amplitude changes:
\(y = A(t)\sin(1000000 \cdot 2\pi t)\)

If \(A(t)\) varies between 0.5 and 1.5, what are the minimum and maximum
peak heights of the signal?

\begin{center}\rule{0.5\linewidth}{0.5pt}\end{center}

\hypertarget{part-c-period-horizontal-stretchcompression-18-problems}{%
\subsection{Part C: Period --- Horizontal Stretch/Compression (18
problems)}\label{part-c-period-horizontal-stretchcompression-18-problems}}

\hypertarget{the-general-form-1}{%
\subsubsection{The General Form}\label{the-general-form-1}}

\[y = \sin(Bx) \quad \text{or} \quad y = \cos(Bx)\]

\textbf{Period} \(= \frac{2\pi}{|B|}\)

\begin{itemize}
\tightlist
\item
  If \(B > 1\): horizontal compression (faster oscillation, shorter
  period)
\item
  If \(0 < B < 1\): horizontal stretch (slower oscillation, longer
  period)
\end{itemize}

\textbf{Frequency} \(= \frac{|B|}{2\pi} = \frac{1}{\text{Period}}\)
(cycles per unit)

\begin{center}\rule{0.5\linewidth}{0.5pt}\end{center}

\textbf{Find the period:}

\textbf{28.} \(y = \sin(2x)\)

\textbf{29.} \(y = \cos(3x)\)

\textbf{30.} \(y = \sin\left(\frac{x}{2}\right)\)

\textbf{31.} \(y = \cos(4\pi x)\)

\textbf{32.} \(y = \sin\left(\frac{\pi x}{3}\right)\)

\begin{center}\rule{0.5\linewidth}{0.5pt}\end{center}

\textbf{Find B if the period is:}

\textbf{33.} Period \(= \pi\)

\textbf{34.} Period \(= 4\pi\)

\textbf{35.} Period \(= 1\) (useful for time in seconds!)

\textbf{36.} Period \(= \frac{1}{60}\) (one cycle per minute)

\begin{center}\rule{0.5\linewidth}{0.5pt}\end{center}

** Science Connection: Musical Frequencies**

\textbf{37.} Middle C has frequency 261.6 Hz. - (a) What is the period
in seconds? - (b) Write the equation for this sound wave:
\(y = \sin(Bt)\) where the period is in seconds.

\textbf{38.} An octave higher means double the frequency. What is the
frequency of C one octave above middle C?

\textbf{39.} The note E above middle C has frequency 329.6 Hz. - (a)
Write the equation. - (b) If you play both C and E together, what
mathematical operation combines them?

\begin{center}\rule{0.5\linewidth}{0.5pt}\end{center}

** Science Connection: Radio Frequencies**

\textbf{40.} FM radio station 98.5 FM broadcasts at 98.5 MHz (megahertz
= million Hz). - (a) What is the period of this electromagnetic wave? -
(b) How does this compare to AM radio at 1 MHz?

\textbf{41.} Wi-Fi operates at either 2.4 GHz or 5 GHz (gigahertz =
billion Hz). - (a) What is the period of a 2.4 GHz wave? - (b) What is
the period of a 5 GHz wave? - (c) Which has more cycles per second?

\begin{center}\rule{0.5\linewidth}{0.5pt}\end{center}

** Science Connection: X-rays and Visible Light**

\textbf{42.} Visible light has frequencies around \(5 \times 10^{14}\)
Hz (500 THz). - (a) What is the period? - (b) Light travels at
\(c = 3 \times 10^8\) m/s. Wavelength \(\lambda = c \cdot T\) where
\(T\) is period. Find the wavelength.

\textbf{43.} X-rays have frequencies around \(10^{18}\) Hz. - (a) What
is the period? - (b) What is the wavelength? - (c) Why can X-rays
``see'' smaller features than visible light?

\begin{center}\rule{0.5\linewidth}{0.5pt}\end{center}

** Science Connection: EUV Lithography**

\textbf{44.} Extreme ultraviolet (EUV) light used in chip manufacturing
has wavelength 13.5 nm. - (a) Convert to meters. - (b) Using
\(\lambda = c \cdot T\), find the period. - (c) Find the frequency.

\textbf{45.} The light in EUV lithography can be modeled as
\(E(t) = E_0\sin(2\pi f t)\). With \(f\) from problem 44: - (a) What is
\(B\) in \(E(t) = E_0\sin(Bt)\)? - (b) This is an incredibly large
number! What does this tell you about how fast the electric field
oscillates?

\begin{center}\rule{0.5\linewidth}{0.5pt}\end{center}

\hypertarget{part-d-phase-shift-horizontal-translation-14-problems}{%
\subsection{Part D: Phase Shift --- Horizontal Translation (14
problems)}\label{part-d-phase-shift-horizontal-translation-14-problems}}

\hypertarget{the-general-form-2}{%
\subsubsection{The General Form}\label{the-general-form-2}}

\[y = \sin(x - C) \quad \text{or} \quad y = \cos(x - C)\]

\textbf{Phase shift} \(= C\) (shifts right if \(C > 0\), left if
\(C < 0\))

Note: In \(y = \sin(Bx - C)\), the phase shift is \(\frac{C}{B}\), not
\(C\)!

\begin{center}\rule{0.5\linewidth}{0.5pt}\end{center}

\textbf{Identify the phase shift:}

\textbf{46.} \(y = \sin\left(x - \frac{\pi}{4}\right)\)

\textbf{47.} \(y = \cos\left(x + \frac{\pi}{3}\right)\)

\textbf{48.} \(y = \sin(2x - \pi)\) --- \emph{Careful! Factor out B
first.}

\textbf{49.} \(y = \cos\left(3x + \frac{\pi}{2}\right)\)

\begin{center}\rule{0.5\linewidth}{0.5pt}\end{center}

\textbf{Write the equation:}

\textbf{50.} Sine function shifted right by \(\frac{\pi}{6}\).

\textbf{51.} Cosine function shifted left by \(\frac{\pi}{4}\).

\textbf{52.} Sine function with period \(\pi\) and phase shift
\(\frac{\pi}{4}\) to the right.

\begin{center}\rule{0.5\linewidth}{0.5pt}\end{center}

** Science Connection: AC Electricity**

\textbf{53.} In three-phase electrical power, three sine waves are used,
each shifted by 120° (or \(\frac{2\pi}{3}\) radians):

\[V_1 = V_0\sin(\omega t)\]
\[V_2 = V_0\sin\left(\omega t - \frac{2\pi}{3}\right)\]
\[V_3 = V_0\sin\left(\omega t - \frac{4\pi}{3}\right)\]

At \(t = 0\), calculate \(V_1\), \(V_2\), and \(V_3\) in terms of
\(V_0\).

\textbf{54.} Show that \(V_1 + V_2 + V_3 = 0\) at \(t = 0\).

\emph{This is always true! The three phases always sum to zero, which is
important for power transmission.}

\begin{center}\rule{0.5\linewidth}{0.5pt}\end{center}

** Science Connection: Signal Phase**

\textbf{55.} GPS works by comparing the phase of signals from multiple
satellites. If two satellites send identical sine waves but one signal
travels farther:

Satellite A: \(y_A = \sin(\omega t)\) Satellite B:
\(y_B = \sin(\omega t - \phi)\)

If the phase difference \(\phi = \frac{\pi}{2}\) and the wave has
frequency 1.575 GHz, how much extra time did signal B take?

\emph{Hint: Phase difference \(\phi\) corresponds to time difference
\(\Delta t = \frac{\phi}{\omega} = \frac{\phi}{2\pi f}\)}

\textbf{56.} At the speed of light (\(3 \times 10^8\) m/s), what extra
distance did signal B travel?

\begin{center}\rule{0.5\linewidth}{0.5pt}\end{center}

** Science Connection: Interference in Thin Films**

\textbf{57.} When light reflects off a thin film (like oil on water),
the two reflections (top and bottom of film) can interfere.

If the film thickness causes a phase shift of \(\phi = \pi\), what
happens when the two waves combine?

\textbf{58.} At what phase shift do you get constructive interference
(bright colors)?

\begin{center}\rule{0.5\linewidth}{0.5pt}\end{center}

\hypertarget{part-e-vertical-shift-moving-the-midline-10-problems}{%
\subsection{Part E: Vertical Shift --- Moving the Midline (10
problems)}\label{part-e-vertical-shift-moving-the-midline-10-problems}}

\hypertarget{the-general-form-3}{%
\subsubsection{The General Form}\label{the-general-form-3}}

\[y = \sin(x) + D \quad \text{or} \quad y = \cos(x) + D\]

\textbf{Vertical shift} \(= D\) (midline moves to \(y = D\))

\textbf{Range:} \([D - A, D + A]\) where \(A\) is amplitude

\begin{center}\rule{0.5\linewidth}{0.5pt}\end{center}

\textbf{Identify amplitude, period, phase shift, and vertical shift:}

\textbf{59.} \(y = 2\sin(x) + 3\)

\textbf{60.} \(y = -3\cos(2x) - 1\)

\textbf{61.} \(y = 4\sin\left(x - \frac{\pi}{2}\right) + 5\)

\textbf{62.} \(y = \frac{1}{2}\cos(4x + \pi) - 2\)

\begin{center}\rule{0.5\linewidth}{0.5pt}\end{center}

** Science Connection: Temperature Cycles**

\textbf{63.} Dublin's average daily temperature varies roughly
sinusoidally through the year: - Average: 10°C - Varies by about ±5°C
from the average - Coldest around January 15, warmest around July 15

Let \(t\) be days after January 1.

\begin{itemize}
\item
  \begin{enumerate}
  \def\labelenumi{(\alph{enumi})}
  \tightlist
  \item
    What is the amplitude?
  \end{enumerate}
\item
  \begin{enumerate}
  \def\labelenumi{(\alph{enumi})}
  \setcounter{enumi}{1}
  \tightlist
  \item
    What is the vertical shift (midline)?
  \end{enumerate}
\item
  \begin{enumerate}
  \def\labelenumi{(\alph{enumi})}
  \setcounter{enumi}{2}
  \tightlist
  \item
    What is the period?
  \end{enumerate}
\item
  \begin{enumerate}
  \def\labelenumi{(\alph{enumi})}
  \setcounter{enumi}{3}
  \tightlist
  \item
    Write an equation using cosine: \(T(t) = A\cos(B(t - C)) + D\)
  \end{enumerate}
\end{itemize}

\emph{Hint: Use cosine with the coldest point as the ``start'' of the
cycle.}

\textbf{64.} Using your equation: - (a) What temperature does it predict
for March 15 (\(t = 73\))? - (b) For June 1 (\(t = 151\))?

\begin{center}\rule{0.5\linewidth}{0.5pt}\end{center}

** Science Connection: Tides**

\textbf{65.} Ocean tides follow roughly sinusoidal patterns (actually
more complex due to Sun and Moon interactions).

A simplified model: High tide at midnight (12:00 AM) with height 3.2 m,
low tide at 6:15 AM with height 0.8 m.

\begin{itemize}
\item
  \begin{enumerate}
  \def\labelenumi{(\alph{enumi})}
  \tightlist
  \item
    What is the amplitude?
  \end{enumerate}
\item
  \begin{enumerate}
  \def\labelenumi{(\alph{enumi})}
  \setcounter{enumi}{1}
  \tightlist
  \item
    What is the midline (average tide level)?
  \end{enumerate}
\item
  \begin{enumerate}
  \def\labelenumi{(\alph{enumi})}
  \setcounter{enumi}{2}
  \tightlist
  \item
    What is the period?
  \end{enumerate}
\item
  \begin{enumerate}
  \def\labelenumi{(\alph{enumi})}
  \setcounter{enumi}{3}
  \tightlist
  \item
    Write the equation using cosine.
  \end{enumerate}
\end{itemize}

\begin{center}\rule{0.5\linewidth}{0.5pt}\end{center}

** Science Connection: Circadian Rhythms**

\textbf{66.} Your body temperature varies throughout the day: - Lowest
around 4:00 AM: 36.2°C - Highest around 4:00 PM: 37.4°C

\begin{itemize}
\item
  \begin{enumerate}
  \def\labelenumi{(\alph{enumi})}
  \tightlist
  \item
    What is the amplitude and midline?
  \end{enumerate}
\item
  \begin{enumerate}
  \def\labelenumi{(\alph{enumi})}
  \setcounter{enumi}{1}
  \tightlist
  \item
    Write an equation where \(t\) is hours after midnight.
  \end{enumerate}
\item
  \begin{enumerate}
  \def\labelenumi{(\alph{enumi})}
  \setcounter{enumi}{2}
  \tightlist
  \item
    What is your predicted body temperature at 8:00 AM?
  \end{enumerate}
\end{itemize}

\begin{center}\rule{0.5\linewidth}{0.5pt}\end{center}

\hypertarget{part-f-the-complete-general-form-10-problems}{%
\subsection{Part F: The Complete General Form (10
problems)}\label{part-f-the-complete-general-form-10-problems}}

\hypertarget{putting-it-all-together}{%
\subsubsection{Putting It All Together}\label{putting-it-all-together}}

\[y = A\sin(B(x - C)) + D \quad \text{or} \quad y = A\cos(B(x - C)) + D\]

\begin{longtable}[]{@{}ll@{}}
\toprule\noalign{}
Parameter & Effect \\
\midrule\noalign{}
\endhead
\bottomrule\noalign{}
\endlastfoot
\(\|A\|\) & Amplitude (vertical stretch) \\
\(A < 0\) & Reflection across midline \\
\(\frac{2\pi}{\|B\|}\) & Period (horizontal stretch) \\
\(C\) & Phase shift (horizontal translation) \\
\(D\) & Vertical shift (midline) \\
\end{longtable}

\begin{center}\rule{0.5\linewidth}{0.5pt}\end{center}

\textbf{For each function, identify all parameters and sketch:}

\textbf{67.} \(y = 3\sin(2x - \pi) + 1\)

\textbf{68.} \(y = -2\cos\left(\frac{x}{2} + \frac{\pi}{4}\right) - 3\)

\begin{center}\rule{0.5\linewidth}{0.5pt}\end{center}

\textbf{Write the equation given:}

\textbf{69.} Amplitude 5, period \(4\pi\), phase shift \(\frac{\pi}{2}\)
right, midline \(y = -2\), using sine.

\textbf{70.} Amplitude 2, period 1 (useful for signals!), no phase
shift, midline \(y = 3\), using cosine.

\begin{center}\rule{0.5\linewidth}{0.5pt}\end{center}

** Science Connection: Musical Synthesis**

\textbf{71.} A synthesizer creates complex sounds by adding sine waves.
A simple ``flute-like'' tone might be:

\[y(t) = 1.0\sin(2\pi \cdot 440t) + 0.3\sin(2\pi \cdot 880t) + 0.1\sin(2\pi \cdot 1320t)\]

\begin{itemize}
\item
  \begin{enumerate}
  \def\labelenumi{(\alph{enumi})}
  \tightlist
  \item
    What are the three frequencies present?
  \end{enumerate}
\item
  \begin{enumerate}
  \def\labelenumi{(\alph{enumi})}
  \setcounter{enumi}{1}
  \tightlist
  \item
    What are the three amplitudes?
  \end{enumerate}
\item
  \begin{enumerate}
  \def\labelenumi{(\alph{enumi})}
  \setcounter{enumi}{2}
  \tightlist
  \item
    The first frequency is the ``fundamental.'' What is the relationship
    between the other frequencies and the fundamental?
  \end{enumerate}
\end{itemize}

\textbf{72.} A square wave can be approximated by:

\[y(t) = \sin(t) + \frac{1}{3}\sin(3t) + \frac{1}{5}\sin(5t) + \frac{1}{7}\sin(7t) + \cdots\]

\begin{itemize}
\item
  \begin{enumerate}
  \def\labelenumi{(\alph{enumi})}
  \tightlist
  \item
    What pattern do you see in the frequencies?
  \end{enumerate}
\item
  \begin{enumerate}
  \def\labelenumi{(\alph{enumi})}
  \setcounter{enumi}{1}
  \tightlist
  \item
    What pattern do you see in the amplitudes?
  \end{enumerate}
\item
  \begin{enumerate}
  \def\labelenumi{(\alph{enumi})}
  \setcounter{enumi}{2}
  \tightlist
  \item
    This is a \textbf{Fourier series}! It shows any periodic wave can be
    built from sine waves.
  \end{enumerate}
\end{itemize}

\begin{center}\rule{0.5\linewidth}{0.5pt}\end{center}

** Science Connection: Quantum Mechanics Preview**

\textbf{73.} In quantum mechanics, a particle in a box has wave
functions that are sine waves:

\[\psi_n(x) = A\sin\left(\frac{n\pi x}{L}\right)\]

where \(L\) is the box length and \(n = 1, 2, 3, ...\)

\begin{itemize}
\item
  \begin{enumerate}
  \def\labelenumi{(\alph{enumi})}
  \tightlist
  \item
    What is the ``period'' of \(\psi_1(x)\)?
  \end{enumerate}
\item
  \begin{enumerate}
  \def\labelenumi{(\alph{enumi})}
  \setcounter{enumi}{1}
  \tightlist
  \item
    What is the ``period'' of \(\psi_2(x)\)?
  \end{enumerate}
\item
  \begin{enumerate}
  \def\labelenumi{(\alph{enumi})}
  \setcounter{enumi}{2}
  \tightlist
  \item
    How many half-wavelengths fit in the box for \(\psi_n(x)\)?
  \end{enumerate}
\end{itemize}

\begin{center}\rule{0.5\linewidth}{0.5pt}\end{center}

** Science Connection: Sampling and Aliasing**

\textbf{74.} When digitizing a signal, you must sample at least twice
per period (Nyquist theorem).

A signal has frequency 1000 Hz. - (a) What is the minimum sampling rate
needed? - (b) If you sample at 800 Hz (too slow), the signal ``aliases''
to a lower frequency. The aliased frequency is \(|f - f_s|\) where
\(f_s\) is the sample rate. What frequency would you hear?

\textbf{75.} CDs sample audio at 44,100 Hz. - (a) What is the highest
frequency that can be accurately recorded? - (b) Why is this sufficient
for human hearing (which goes up to about 20,000 Hz)?

\textbf{76.} A scientist digitizes a signal at 100 samples per second.
She observes a frequency of 30 Hz in the data. Could the original signal
actually have been 70 Hz? Explain.

\begin{center}\rule{0.5\linewidth}{0.5pt}\end{center}

\end{document}
